\chapter{Implementation in C++}
\label{c:implementation}

%%%%%%%%%%%%%%%%%%%%%%%%%%%%%%%%%%%%%%%%%%%%%%%%%%%%%%%%%%%%%%%%%%%%%%%%%%%%%%%%
% ==========================================
% Section: Core Data Structures
% ==========================================
\section{Core Data Structures}
To construct the foundation of the simulator, we must first define the core data structures that map the mathematical framework into efficient computational memory.

% ==========================================
% Section: Defining the state
% ==========================================
\subsection{Defining the state}
As established in Sections~\ref{subsec:quantum-states} and~\ref{subsec:tensor-products}, a quantum state is represented as a linear combination of orthogonal basis states $\ket{i}$ and their corresponding complex probability amplitudes $c_i$. Computationally, simulating these many-body states requires a data structure optimized for rapid access and modification. This is particularly crucial because, during generator-based evolution, the algorithm must dynamically retrieve and manipulate specific pairs of interacting basis states.

To represent the fermionic Fock states, we use the occupation number representation defined in Equation \eqref{eq:occupation_number_representation}. Because fermions obey the Pauli exclusion principle, a given orbital is either strictly empty or occupied. This binary nature maps perfectly to a bitstring. In C++, this is naturally and efficiently encoded as an unsigned 64-bit integer (\texttt{uint64\_t}). Utilizing native integer types allows the simulator to execute optimized, hardware-level bitwise operations when applying creation and annihilation operators. For the associated probability amplitudes, we utilize the standard library's double-precision complex type, \texttt{std::complex<double>}, ensuring robust numerical stability during phase accumulations.

To bind the basis states to their amplitudes, we require a structure that supports dynamic sizing and fast lookups. Rather than relying on dense arrays—which would scale exponentially and store unnecessary zero-amplitude states—we implement a sparse representation using a hash map. Specifically, we utilize the \texttt{ankerl::unordered\_dense::map}. Unlike the standard \texttt{std::unordered\_map}, this optimized library provides a cache-friendly, contiguous memory layout.

Combining these architectural choices, the complete quantum state is elegantly defined via a single type alias:

\begin{minted}{cpp}
    using State = ankerl::unordered_dense::map<uint64_t, std::complex<double>>;
\end{minted}

% ==========================================
% Section: Excitations
% ==========================================
\section{Excitations}
Building upon the core state representation, we now address the application of excitation operators. We first examine the simplified qubit excitation baseline before introducing the necessary fermionic modifications.

% ==========================================
% Subsection: Qubit Baseline
% ==========================================
\subsection{Qubit Baseline}
As discussed in Section~\ref{subsec:qubit-baseline}, implementing the qubit excitation requires three key parameters: the initial quantum state, the orbital sets $\boldsymbol{k}$ and $\boldsymbol{l}$ (which define the electron transitions), and the rotation angle $\theta$.

To avoid unnecessary memory allocations, the orbital sets are passed as constant references to standard integer vectors (\texttt{std::vector<int>}), while the angle $\theta$ is provided as a double-precision float. This yields the following function signature:

\begin{minted}{cpp}
State apply_qubit_excitation(const State& state, const std::vector<int>& k, const std::vector<int>& l, const double theta)
\end{minted}

Memory for the \texttt{new\_state} map is preallocated to establish a baseline capacity, minimizing dynamic reallocations when inserting known basis states and their phase factors. If an excitation generates a previously unseen partner state that exceeds this capacity, the underlying C++ container dynamically allocates the required memory. The core algorithm then iterates over all non-zero basis states, applying the unitary transformation defined in Equation \eqref{eq:qubit-excitation-unitary-transformation-configuration-differences}:

\begin{equation} \tag{\ref{eq:qubit-excitation-unitary-transformation-configuration-differences}}
    \qubitExcitationUnitaryTransformationConfigurationDifferences
\end{equation}

To iterate efficiently over the sparse state map, the implementation uses structured bindings. This allows direct, reference-based extraction of the basis state and its corresponding amplitude without invoking costly memory copies:

\begin{minted}{cpp}
    for (const auto& [basis_state, coeff] : state)
\end{minted}

Before performing floating-point arithmetic or allocating new state entries, the simulator evaluates the nullspace projector $P_0$. A basis state belongs to the active space if and only if all origin orbitals $\boldsymbol{k}$ are uniformly occupied while all target orbitals $\boldsymbol{l}$ are uniformly unoccupied, or vice versa.

Because the basis state is encoded as a \texttt{uint64\_t}, this physical condition is evaluated using raw hardware-level bitwise operations. Bit shifting (\texttt{>>}) and masking (\texttt{\& 1ULL}) execute in a single CPU cycle. Looping over the orbital indices allows the algorithm to skip nullspace states with minimal computational overhead:

\begin{minted}{cpp}
    bool k_val = (basis_state >> k[0]) & 1ULL;
    bool skip = false;
    for (size_t i = 0; i < k.size(); i++) {
        if (((basis_state >> k[i]) & 1ULL) != k_val) skip = true;
        if (((basis_state >> l[i]) & 1ULL) == k_val) skip = true;
    }
\end{minted}

States in the nullspace undergo no rotation or excitation, but they must still be copied into \texttt{new\_state} to preserve the unaffected components of the wavefunction.

For basis states in the active space, the simulator applies the excitation. First, the phase of the original basis state is updated according to Equation \eqref{eq:qubit-excitation-update-basis-state}:

\begin{equation} \tag{\eqref{eq:qubit-excitation-update-basis-state}}
    c_{\boldsymbol{i}} \leftarrow c_{\boldsymbol{i}} \cos\left(\frac{\theta}{2}\right)
\end{equation}

In C++, this calculation is implemented as:

\begin{minted}{cpp}
    std::complex<double> c_ibs = coeff * std::cos(theta / 2);
\end{minted}

Next, the partner basis state is calculated using Equations \eqref{eq:qubit-excitation-baseline-partner-state-1} and \eqref{eq:qubit-excitation-baseline-partner-state-2}:

\begin{align}
    c \ket{\boldsymbol{i'}} &= -iG^{\boldsymbol{k}}_{\boldsymbol{l}} \ket{\boldsymbol{i}}, \tag{\eqref{eq:qubit-excitation-baseline-partner-state-1}} \\
    -iG^{\boldsymbol{k}}_{\boldsymbol{l}} &= \left( \prod_n \sigma^-_{k_n} \sigma^+_{l_n} - \text{h.c.} \right). \tag{\eqref{eq:qubit-excitation-baseline-partner-state-2}}
\end{align}

Computationally, the electron transitions correspond to toggling the bits at the positions defined by the orbital sets $\boldsymbol{k}$ and $\boldsymbol{l}$. Since the state is already confirmed to be in the active space, applying a bitwise XOR (\texttt{\^{}}) correctly removes and adds electrons at the respective orbitals:

\begin{minted}{cpp}
    uint64_t partner_basis_state = basis_state;
    for (int i = 0; i < k.size(); i++) {
        partner_basis_state ^= (1ULL << k[i]);
        partner_basis_state ^= (1ULL << l[i]);
    }
\end{minted}

The direction of the excitation, stored in $c$, is determined by checking the occupancy of the first orbital in the $\boldsymbol{k}$ set:

\begin{minted}{cpp}
    int c = (basis_state >> k[0]) & 1ULL ? 1 : -1;
\end{minted}

Finally, the amplitude of the newly acquired partner state is calculated using Equation \eqref{eq:qubit-excitation-partner-state-update-phase}:

\begin{equation} \tag{\ref{eq:qubit-excitation-partner-state-update-phase}}
    c_{\boldsymbol{i'}} \leftarrow c_{\boldsymbol{i'}} + c \sin\left(\frac{\theta}{2}\right).
\end{equation}

This maps to the following code snippet:

\begin{minted}{cpp}
    std::complex<double> c_ipbs = c * std::sin(theta/2) * coeff;
\end{minted}

Combining these components yields the complete qubit excitation function:

\begin{minted}{cpp}
    State apply_qubit_excitation(const State& state, const std::vector<int>& k, const std::vector<int>& l, const double theta) {
    // Initials
    State new_state;
    new_state.reserve(state.size() * 2);
    // Loop through state
    for (const auto& [basis_state, coeff] : state) {
        // Check for P0
        bool k_val = (basis_state >> k[0]) & 1ULL;
        bool skip = false;
        for (int i = 0; i < k.size(); i++) {
            if (((basis_state >> k[i]) & 1ULL) != k_val) skip = true;
            if (((basis_state >> l[i]) & 1ULL) == k_val) skip = true;
        }
        if (skip) {
            // Nullspace
            new_state[basis_state] += coeff;
        } else {
            // Active Space - Implement excitation
            std::complex<double> c_ibs = coeff * std::cos(theta / 2);
            uint64_t partner_basis_state = basis_state;
            for (int i = 0; i < k.size(); i++) {
                partner_basis_state ^= (1ULL << k[i]);
                partner_basis_state ^= (1ULL << l[i]);
            }
            int c = (basis_state >> k[0]) & 1ULL ? 1 : -1;
            std::complex<double> c_ipbs = c * std::sin(theta/2) * coeff;

            new_state[basis_state] += c_ibs;
            new_state[partner_basis_state] += c_ipbs;
        }
    }
    return new_state;
}
\end{minted}

% ==========================================
% Subsection: introducing fermions
% ==========================================
\subsection{introducing fermions}
After having the Baseline of the Qubit defined we can now concentrate on the differences of the fermionic Excitation. As also already seen before in Section \eqref{sec:matrix-free-direct-state-evolution} the changes are not severe. Actually we only need to apply more logic to the Excitation itself because the structure is the same and only the Generator of the unitary funtion changed.

The initialization and the checking for nullspace is exactly the same as before, but the computation of the excitation changes. First we need to calculate the direction in which the operators act. For that we take the boolean value (\texttt{k\_val}) of our first orbital in the orbital set \texttt{k}, which we previously already calculated in the nullspace check as following:
\begin{minted}{cpp}
    bool k_val = (basis_state >> k[0]) & 1ULL;
\end{minted}
We can then check if this \texttt{k\_val} is empty or occupied, which tells us the direction. Empty means that \texttt{k} is the creation set, so the set of orbitals where an electron must be created. Therefore \texttt{l} is the annihilation set, so the set of orbitals which are all occupied and the electrons needed to be annihilated. Note that we can only look at the first because we already applied the nullspace Projector and therefore know that this basis state is in the active space. In C++ we can use two ternary operators to describes this behaviour as following:
\begin{minted}{cpp}
    const std::vector<int>& annihilate_idx = (k_val == 0) ? l : k;
    const std::vector<int>& create_idx = (k_val == 0) ? k : l;
\end{minted}

Now we can apply the excitation to a copy of the state (\texttt{new\_basis\_state}), because we later need the original one again to update the amplitudes correctly. As explained before in \eqref{subsec:incorporating-the-fermionic-phase}, we need to define the behavior of equation \eqref{eq:annhiliation-creation-algebraic-definition}:
\begin{align} \tag{\eqref{eq:annhiliation-creation-algebraic-definition}}
    a_k \ket{\mathbf{n}} &= \delta_{n_k, 1} (-1)^{\sum_{j=0}^{k-1} n_j} \ket{n_0, \dots, n_k - 1, \dots, n_{M-1}}, \\
    a_k^\dagger \ket{\mathbf{n}} &= \delta_{n_k, 0} (-1)^{\sum_{j=0}^{k-1} n_j} \ket{n_0, \dots, n_k + 1, \dots, n_{M-1}}.
\end{align}

We start with the phase factor $(-1)^{\sum_{j=0}^{k-1} n_j}$. To translate it into C++ we build a mask, which has 1 from the first index to then the current index, like this:
\begin{minted}{cpp}
    uint64_t mask = (1ULL << idx) - 1;
\end{minted}

After this we can count the number of occupied orbitals using the mask. For that we apply the mask with a $\&$ operator to the basis state. That cuts out the part of the basis state we need for the computation. On this part we then can apply the C++ \texttt{\_\_builtin\_popcountll} function which counts the number of set bits in the bitstring. After that we only need to check if this number is even or odd, which we can also do with a bit operation \texttt{\% 2}. If this result is now not 0 it means its uneven and we need to apply a phase of -1. After this we only need to apply the bitflip of the orbital to annihilate or create the electron. The code for this looks like this:
\begin{minted}{cpp}
    for (int idx : annihilate_idx) {
        uint64_t mask = (1ULL << idx) - 1;
        if (__builtin_popcountll(new_basis_state & mask) % 2 != 0) {
            phase = -phase;
        }
        new_basis_state ^= (1ULL << idx);
    }
\end{minted}

Finally we only need to update the amplitudes correctly, for that we first need to look at the direction of the excitation to define the \texttt{c} value correctly. We can do this in the same manner we did it before for the orbital sets:
\begin{minted}{cpp}
    int c = (k_val == 1) ? -1 : 1;
\end{minted}
The rest of the amplitudes get updated excatly as we alredy did in the qubit implementation.

Putting all toegther we get:
\begin{minted}{cpp}
    State apply_fermion_excitation(const State& state, const std::vector<int>& k,
    const std::vector<int>& l, const double theta) {
    // Checks

    // Initials
    State new_state;
    new_state.reserve(state.size() * 2);
    // Loop through state
    for (const auto& [basis_state, coeff] : state) {
        // Check for P0
        bool k_val = (basis_state >> k[0]) & 1ULL;
        bool skip = false;
        for (int i = 0; i < k.size(); i++) {
            if (((basis_state >> k[i]) & 1ULL) != k_val) skip = true;
            if (((basis_state >> l[i]) & 1ULL) == k_val) skip = true;
        }
        if (skip) {
            new_state[basis_state] += coeff;
        } else {
            // Compute the phase
            int phase = 1;
            uint64_t new_basis_state = basis_state;

            // Determine operation direction:
            // If k_val == 0, 'k' is empty and 'l' is occupied. (Applying T)
            // If k_val == 1, 'k' is occupied and 'l' is empty. (Applying T^\dagger)
            const std::vector<int>& annihilate_idx = (k_val == 0) ? l : k;
            const std::vector<int>& create_idx = (k_val == 0) ? k : l;

            // Apply annihilation
            for (int idx : annihilate_idx) {
                uint64_t mask = (1ULL << idx) - 1;
                if (__builtin_popcountll(new_basis_state & mask) % 2 != 0) {
                    phase = -phase;
                }
                new_basis_state ^= (1ULL << idx);
            }

            // Apply creation
            for (int idx : create_idx) {
                uint64_t mask = (1ULL << idx) - 1;
                if (__builtin_popcountll(new_basis_state & mask) % 2 != 0) {
                    phase = -phase;
                }
                new_basis_state ^= (1ULL << idx);
            }

            // Relative sign for T or -T^\dagger
            int c = (k_val == 1) ? -1 : 1;

            // update amplitudes
            new_state[basis_state] += coeff * std::cos(theta / 2);
            new_state[new_basis_state] += c * phase * std::sin(theta / 2) * coeff;
        }
    }
    return new_state;
}
\end{minted}

% ==========================================
% Subsection: Fermionic Excitation Implementation
% ==========================================
\subsection{Fermionic Excitation Implementation}\label{subsec:fermionic-excitation-implementation}
Building upon the matrix-free evolution framework from Section \ref{sec:matrix-free-direct-state-evolution}, implementing fermionic excitations requires incorporating anticommutation phase logic. While the structural state evolution remains identical to the qubit baseline, the generator action must account for fermionic statistics.

Following the nullspace check, the occupation of the first target orbital (\texttt{k\_val}) determines the direction of the operator (i.e., whether $T$ or $T^\dagger$ is applied):
\begin{minted}{cpp}
    bool k_val = (basis_state >> k[0]) & 1ULL;
\end{minted}
Because the state is guaranteed to be within the active space by the nullspace projector, this single bit determines the operation roles: if \texttt{k\_val} is 0, set \bold{k} corresponds to creation operators and set \bold{l} to annihilation operators. Conversely, if \texttt{k\_val} is 1, the roles reverse. In C++, this routing is handled via references to avoid unnecessary memory copying:
\begin{minted}{cpp}
    const std::vector<int>& annihilate_idx = (k_val == 0) ? l : k;
    const std::vector<int>& create_idx = (k_val == 0) ? k : l;
\end{minted}

To evaluate the algebraic action defined in Eq. \ref{eq:annhiliation-creation-algebraic-definition}, the excitation is applied to a working copy of the bitstring (\texttt{new\_basis\_state}).
\begin{align} \tag{\ref{eq:annhiliation-creation-algebraic-definition}}
    a_k \ket{\mathbf{n}} &= \delta_{n_k, 1} (-1)^{\sum_{j=0}^{k-1} n_j} \ket{n_0, \dots, n_k - 1, \dots, n_{M-1}}, \\
    a_k^\dagger \ket{\mathbf{n}} &= \delta_{n_k, 0} (-1)^{\sum_{j=0}^{k-1} n_j} \ket{n_0, \dots, n_k + 1, \dots, n_{M-1}}.
\end{align}
Evaluating the parity phase $(-1)^{\sum_{j=0}^{k-1} n_j}$ requires counting the set bits preceding index \texttt{idx}. This is achieved efficiently by generating a bitmask, isolating the preceding orbitals, and invoking the hardware-accelerated population count instruction (\texttt{\_\_builtin\_popcountll}). If the resulting count is odd, the global phase is inverted. The orbital state is subsequently flipped using a bitwise XOR operation:
\begin{minted}{cpp}
    for (int idx : annihilate_idx) {
        uint64_t mask = (1ULL << idx) - 1;
        if (__builtin_popcountll(new_basis_state & mask) % 2 != 0) {
            phase = -phase;
        }
        new_basis_state ^= (1ULL << idx);
    }
\end{minted}

Finally, the amplitudes are updated. The relative sign factor $c$ is determined by the excitation direction, and the trigonometric coefficients are applied to map the transition:
\begin{minted}{cpp}
    int c = (k_val == 1) ? -1 : 1;
\end{minted}

The complete function integrating the nullspace projection, phase tracking, and amplitude updates is implemented as follows:
\begin{minted}{cpp}
State apply_fermion_excitation(const State& state, const std::vector<int>& k,
                               const std::vector<int>& l, const double theta) {
    State new_state;
    new_state.reserve(state.size() * 2);

    for (const auto& [basis_state, coeff] : state) {
        // Nullspace projector (P0 check)
        bool k_val = (basis_state >> k[0]) & 1ULL;
        bool skip = false;
        for (size_t i = 0; i < k.size(); i++) {
            if (((basis_state >> k[i]) & 1ULL) != k_val) skip = true;
            if (((basis_state >> l[i]) & 1ULL) == k_val) skip = true;
        }

        if (skip) {
            new_state[basis_state] += coeff;
        } else {
            int phase = 1;
            uint64_t new_basis_state = basis_state;

            // Determine operator direction (T vs T^\dagger)
            const std::vector<int>& annihilate_idx = (k_val == 0) ? l : k;
            const std::vector<int>& create_idx = (k_val == 0) ? k : l;

            // Apply annihilation operators sequentially
            for (int idx : annihilate_idx) {
                uint64_t mask = (1ULL << idx) - 1;
                if (__builtin_popcountll(new_basis_state & mask) % 2 != 0) {
                    phase = -phase;
                }
                new_basis_state ^= (1ULL << idx);
            }

            // Apply creation operators sequentially
            for (int idx : create_idx) {
                uint64_t mask = (1ULL << idx) - 1;
                if (__builtin_popcountll(new_basis_state & mask) % 2 != 0) {
                    phase = -phase;
                }
                new_basis_state ^= (1ULL << idx);
            }

            // Relative sign for T or -T^\dagger
            int c = (k_val == 1) ? -1 : 1;

            // Update amplitudes
            new_state[basis_state] += coeff * std::cos(theta / 2);
            new_state[new_basis_state] += c * phase * std::sin(theta / 2) * coeff;
        }
    }
    return new_state;
}
\end{minted}

% ==========================================
% Section: fSWAP
% ==========================================
\section{fSWAP Implementation}
To avoid constructing the dense operator matrix, the simulator evaluates the fSWAP gate directly on the state vector. Iterating exclusively over the non-zero amplitudes reduces time and memory complexity. Bitwise shifts and XOR masks manipulate orbital occupancies directly, maximizing execution speed and avoiding unnecessary memory allocations:

\begin{minted}{cpp}
const uint64_t mask_i = 1ULL << i;
const uint64_t mask_j = 1ULL << j;

for (const auto& [basis_state, coeff] : state) {
    const bool bit_i = (basis_state >> i) & 1ULL;
    const bool bit_j = (basis_state >> j) & 1ULL;

    if (bit_i == bit_j) {
        // |00⟩ → |00⟩, |11⟩ → -|11⟩
        const double phase = bit_i ? -1.0 : 1.0;
        new_state[basis_state] += phase * coeff;
    } else {
        // |01⟩ ↔ |10⟩: apply XOR bit-swap
        const uint64_t swapped = basis_state ^ mask_i ^ mask_j;
        new_state[swapped] += coeff;
    }
}
\end{minted}


% ==========================================
% Section: Fermionic Excpectation Value
% ==========================================
\section{Fermionic Expectation Value}
Following the mathematical decomposition in Section~\ref{sec:expectation-value-evaluation}, the simulator evaluates the expectation value $\bra{\phi} H \ket{\psi}$ by accumulating the independent contributions of each Hamiltonian term:

\begin{minted}{cpp}
for (const auto& term : operator_sum) {
    result += expectation_value_fermionic_term(phi, psi, term);
}
\end{minted}

For each fermionic operator string $\hat{O}_k$, \texttt{expectation\_value\_fermionic\_term} iterates exclusively over the non-zero amplitudes of $\ket{\psi}$. Annihilation and creation operators are applied in the same way as already before in chapter~\ref{subsec:fermionic-excitation-implementation}

To evaluate the inner product without full matrix multiplication, the simulator projects the transformed basis state $\ket{x'_k}$ onto $\bra{\phi}$. Utilizing a hash-table lookup (\texttt{find()}) to retrieve the conjugate amplitude $d_{x'_k}^*$, directly implementing the collapsed summation:
\begin{equation}
    \bra{\phi} H \ket{\psi} = \sum_{k \in H} \sum_{x \in \psi} \underbrace{d_{x'_k}^*}_{\text{coeff\_phi}} \cdot \underbrace{w_k}_{\text{weight}} \cdot \underbrace{p_{k,x}}_{\text{phase}} \cdot \underbrace{c_x}_{\text{coeff\_psi}}
\end{equation}

\begin{minted}{cpp}
auto it = phi.find(new_basis_state);
if (it != phi.end()) {
    const std::complex<double> coeff_phi = std::conj(it->second);
    term_overlap += coeff_phi * (term.weight * std::complex<double>(phase, 0.0) * coeff_psi);
}
\end{minted}





